\chapter{Objetivos}

Neste capítulo são apresentados o objetivo geral e os objetivos específicos que norteiam o desenvolvimento deste trabalho de conclusão de curso.

\section{Objetivo Geral}

Dado o contexto da pesquisa, o objetivo geral deste trabalho é desenvolver, implementar e avaliar modelos baseados em aprendizado contrastivo para a predição de propriedades físico-químicas de materiais bidimensionais, buscando gerar representações latentes informativas que superem as limitações de modelos puramente supervisionados em cenários de escassez de dados rotulados.

\section{Objetivos Específicos}

Para atingir o objetivo geral, foram definidos os seguintes objetivos específicos:

\begin{itemize}
    \item Realizar o levantamento bibliográfico sobre aprendizado contrastivo (auto-supervisionado e supervisionado) e  representações estruturais aplicadas à ciência computacional de materiais;
    \item Estruturar e pré-processar conjuntos de dados de materiais bidimensionais (como a base C2DB), extraindo descritores estruturais e representações adequadas à modelagem;
    \item Implementar e treinar arquiteturas de aprendizado contrastivo (incluindo redes siamesas e codificadores baseados em métricas de similaridade) para o mapeamento dos materiais em um espaço latente de baixa dimensão;
    \item Avaliar a qualidade das representações latentes aprendidas e o desempenho preditivo para propriedades como \textit{band gap} e energia de formação, comparando os resultados com métodos tradicionais previamente explorados.
\end{itemize}
